\documentclass[11pt,letterpaper]{article}

\usepackage[T1]{fontenc}
\usepackage[utf8]{inputenc}
\usepackage{lmodern}
\usepackage[letterpaper,margin=0.94in]{geometry}
\usepackage{microtype}
\usepackage{amsmath,amssymb,amsthm,mathtools}
\usepackage{mathrsfs}
\usepackage{enumitem}
\usepackage{xcolor}
\usepackage{needspace}
\usepackage{hyperref}
\usepackage{cite}

\hypersetup{
  colorlinks=true,
  linkcolor=black,
  citecolor=black,
  urlcolor=black,
  pdftitle={Global Regularity for the Three-Dimensional Incompressible Navier–Stokes Equations},
  pdfauthor={Thomas Birnie},
  pdfsubject={The datum-generated compatible intersection as an estimate generator},
  pdfkeywords={Navier–Stokes, global regularity, compatible intersection, mixed jet}
}

\setlength{\parindent}{1.25em}
\setlength{\parskip}{0.16em}
\setlength{\emergencystretch}{2em}
\setlist[itemize]{leftmargin=1.5em,itemsep=0.18em,topsep=0.3em}

\newtheorem{theorem}{Theorem}
\newtheorem{lemma}[theorem]{Lemma}
\newtheorem{proposition}[theorem]{Proposition}
\newtheorem{corollary}[theorem]{Corollary}
\theoremstyle{definition}
\newtheorem{definition}[theorem]{Definition}
\theoremstyle{remark}
\newtheorem{remark}[theorem]{Remark}

\newcommand{\R}{\mathbb R}
\newcommand{\Nzero}{\mathbb N_0}
\newcommand{\Hs}{H^\infty}
\newcommand{\cR}{\mathcal R}
\newcommand{\cI}{\mathcal I}
\newcommand{\frakR}{\mathfrak R}
\newcommand{\Pdiv}{\mathbb P\nabla\!\cdot}
\newcommand{\norm}[1]{\left\lVert #1\right\rVert}
\newcommand{\dd}{\,\mathrm d}
\newcommand{\vort}{\omega}
\newcommand{\strain}{S}
\newcommand{\Istar}{I_*}
\newcommand{\pair}[2]{\left\langle #1,#2\right\rangle}

\title{\vspace{-1.55em}\textbf{Global Regularity for the Three-Dimensional\\
Incompressible Navier{--}Stokes Equations}\\[0.42em]
\large A proof through the datum-generated compatible intersection}
\author{Thomas Birnie}
\date{August 2026}

\begin{document}
\maketitle
\vspace{-1.1em}

\begin{abstract}
For one smooth divergence-free datum and one fixed \(\nu>0\), the
three-dimensional incompressible equation generates one maximal
velocity--pressure history.  The equation also generates the complete mixed
derivative jet of that history.  The datum-generated whole-terminal theorem
places every finite part of this jet in a compatible Sobolev rectangle, and
the rectangles assemble into one projective all-orders intersection.  That
intersection is an estimate generator: every fixed finite-order continuation
bound in the paper is obtained by selecting a sufficiently high coordinate.
The resulting bounds give one common \(H^\infty\) endpoint, the actual
pressure-compatible endpoint jet, and a smooth restart of the same history
across any proposed finite terminal time.  The maximal time is infinite.
\end{abstract}

\section*{Introduction}

Fix one smooth divergence-free finite-energy datum \(u_0\) in three space
dimensions and one parameter \(\nu>0\).  Local theory generates one maximal
velocity--pressure history \((u,p)\) on \([0,T_*)\).  The question is whether
its maximal time can be finite.  Energy controls the velocity itself, but it
does not by itself place every fixed temporal and spatial response in a
bound that survives all the way to \(T_*\).  A finite terminal face could
remain only if some response needed for continuation escaped.  The proof
follows this same history until that possibility is closed.

The datum first generates one maximal velocity--pressure pair.  The equation then
generates its temporal and spatial derivative jet, including the pressure
at each order.  The first theorem of the proof places every finite portion
of that jet in a whole-terminal Sobolev rectangle and makes the rectangles
compatible under restriction.  Their projective intersection retains every
finite temporal and spatial coordinate at once.  Once that intersection has
been formed, the estimates usually pursued one at a time become projections
of a single stronger object.

This order is essential.  The Serrin row, the critical
\(L_t^\infty L_x^3\) row, the vorticity and strain actions, the critical
height, the dyadic shell sums, and the pressure bounds are not used to
produce the intersection.  They are read from it after it exists.  The same
intersection then supplies a common terminal value in every Sobolev
topology, identifies the incoming pressure-complete jet, and permits a local
restart that uniqueness glues to the original history.

The classical background is well established.  Leray constructed global
weak solutions \cite{Leray1934}; Prodi and Serrin identified scale-sensitive
regularity classes \cite{Prodi1959,Serrin1962,Serrin1963}; Fujita--Kato and
Kato developed the strong local theory used at the restart
\cite{FujitaKato1964,Kato1984}; and the critical
\(L_t^\infty L_x^3\) endpoint was resolved by
Escauriaza, Seregin, and \v{S}ver\'ak
\cite{EscauriazaSereginSverak2003}.  The role of those results in this paper
is sharply delimited.  They supply standard local, embedding, or comparison
interfaces.  The proof's central object is the datum-generated compatible
intersection stated in Theorem~\ref{thm:intersection}.

The paper is organized by the causal order of the argument.  Sections
1--2 fix the one history and derive its pressure-complete jet.  Section~3
states the whole-terminal rectangle theorem and proves the projective
assembly.  Section~4 turns the intersection into a general estimate
generator.  Section~5 develops the critical adjacent block and the full
classical consequence suite.  Sections~6--8 construct the common endpoint,
identify its actual jet, and restart the same history.  The appendix records
the finite-order analytic facts used in those passages.

\paragraph{Conventions.}
All velocity Sobolev spaces are vector-valued over \(\R^3\); scalar spaces
are used for pressure.  We write
\[
H^\infty(\R^3):=\bigcap_{m=0}^{\infty}H^m(\R^3),
\qquad
H^\infty_\sigma(\R^3)
:=\{v\in H^\infty(\R^3;\R^3):\nabla\cdot v=0\}.
\]
All time derivatives are distributional derivatives until continuity
selects their canonical representatives.  Repeated spatial indices are
summed.  A \emph{pressure-complete history} means one velocity field,
its pressure recovered from that same field by the fixed normalization
below, and all differentiated velocity and pressure rows generated from the
same equation.

\begin{theorem}[Global regularity]
\label{thm:global}
Let \(\nu>0\), and let
\[
 u_0\in H^\infty(\R^3;\R^3)\cap L^2(\R^3;\R^3),
 \qquad \nabla\cdot u_0=0.
\]
The three-dimensional incompressible Navier{--}Stokes system on \(\R^3\)
has a unique global smooth velocity--pressure history generated by \(u_0\),
with pressure fixed by the decaying Riesz-transform normalization.  For
every \(t\ge0\),
\begin{equation}
\norm{u(t)}_2^2
+2\nu\int_0^t\norm{\nabla u(s)}_2^2\dd s
=\norm{u_0}_2^2.
\label{eq:global-energy}
\end{equation}
In particular,
\[
 u,p\in C^\infty([0,\infty)\times\R^3).
\]
\end{theorem}

\paragraph{Proof of Theorem~\ref{thm:global}.}
The proof follows one object.  Its first movement is
\begin{equation}
\boxed{
u_0
\longrightarrow (u,p)
\longrightarrow \{J_{n,\alpha},\Pi_{n,\alpha}\}_{n,\alpha}
\longrightarrow \{\cR_{N,M}[u_0]\}_{N,M}
\longrightarrow \cI[u_0].}
\label{eq:opening-chain}
\end{equation}
Every arrow in \eqref{eq:opening-chain} retains the datum, the parameter,
the pressure convention, and the history.  The point is not to collect
separate estimates.  It is to construct one object from which each requested
finite estimate can be read.

\section{One datum generates one pressure-complete history}

Standard smooth-data local theory gives a unique classical solution and a
maximal lifespan; see \cite{FujitaKato1964,Kato1984}.  Let \((u,p)\) be that
maximal pressure-complete history generated by \(u_0\), defined on
\([0,T_*)\).  On every strict subinterval \([0,T]\subset[0,T_*)\), it is
smooth and satisfies
\begin{equation}
\partial_tu+(u\cdot\nabla)u+\nabla p=\nu\Delta u,
\qquad
\nabla\cdot u=0,
\qquad
u(0)=u_0.
\label{eq:ns}
\end{equation}
Testing the momentum equation against \(u\), using
\(\nabla\cdot u=0\), and integrating by parts gives
\begin{equation}
\norm{u(t)}_2^2
+2\nu\int_0^t\norm{\nabla u(s)}_2^2\dd s
=\norm{u_0}_2^2
\qquad(0\le t<T_*).
\label{eq:energy-before-endpoint}
\end{equation}
This identity will pass through the restart once the terminal value and the
same-history continuation have been identified.
Write
\[
R_j:=\partial_j(-\Delta)^{-1/2}.
\]
Taking the divergence of \eqref{eq:ns} and fixing the decaying normalization
determines pressure from this same velocity field:
\begin{equation}
-\Delta p=\partial_i\partial_j(u_i u_j),
\qquad
p=R_iR_j(u_i u_j).
\label{eq:pressure}
\end{equation}
The velocity and pressure are not two histories joined after the fact.
Equation \eqref{eq:pressure} records the pressure response carried by the
same evolving velocity field.

Assume for contradiction that
\[
T_*<\infty.
\]
All finite-horizon time estimates below are made on this contradiction
branch.  No later step changes the object named in \eqref{eq:ns}.

\section{The equation generates the complete mixed jet}

Set
\[
V_n:=\partial_t^n u,
\qquad
Q_n:=\partial_t^n p,
\qquad n\in\Nzero.
\]
Repeated differentiation of \eqref{eq:pressure} gives
\begin{equation}
Q_n
=R_iR_j\sum_{a=0}^{n}\binom{n}{a}V_{a,i}V_{n-a,j}.
\label{eq:pressure-recurrence}
\end{equation}
Differentiating \eqref{eq:ns} gives the matching velocity recurrence
\begin{equation}
V_{n+1}
=\nu\Delta V_n
-\sum_{a=0}^{n}\binom{n}{a}(V_a\cdot\nabla)V_{n-a}
-\nabla Q_n.
\label{eq:velocity-recurrence}
\end{equation}
At height \(n\), the recurrence has not introduced another solution.  It has
exposed the next temporal response of the history already generated by
\(u_0\).

For a spatial multi-index \(\alpha\), define
\[
J_{n,\alpha}:=\partial_t^n\partial_x^\alpha u,
\qquad
\Pi_{n,\alpha}:=\partial_t^n\partial_x^\alpha p.
\]
For multi-indices, \(\beta\le\alpha\) means coordinatewise inequality, and
\[
\binom{\alpha}{\beta}
:=\prod_{\ell=1}^{3}\binom{\alpha_\ell}{\beta_\ell}.
\]
Leibniz' rule applied to \eqref{eq:ns} yields the full mixed row
\begin{equation}
\begin{aligned}
J_{n+1,\alpha}
&+\sum_{a=0}^{n}\sum_{\beta\le\alpha}
  \binom{n}{a}\binom{\alpha}{\beta}
  (J_{a,\beta}\cdot\nabla)J_{n-a,\alpha-\beta}
  +\nabla\Pi_{n,\alpha} \\
&=\nu\Delta J_{n,\alpha},
\qquad
\nabla\cdot J_{n,\alpha}=0.
\end{aligned}
\label{eq:mixed-recurrence}
\end{equation}
The pressure coordinate in this row is determined by
\begin{equation}
-\Delta\Pi_{n,\alpha}
=\partial_i\partial_j\partial_t^n\partial_x^\alpha(u_i u_j).
\label{eq:mixed-pressure}
\end{equation}
Equations \eqref{eq:pressure-recurrence}--\eqref{eq:mixed-pressure} form the
pressure-complete mixed jet of the original history.  Its temporal rows,
spatial Sobolev rows, vorticity, strain, critical-height coordinates, and
later continuation quantities are all finite views of this one jet.

\begin{proposition}[The jet is generated by the datum]
\label{prop:datum-jet}
For every \(n\in\Nzero\), the initial values \(V_n(0)\) and \(Q_n(0)\) are
uniquely determined by \(u_0\), \(\nu\), and the pressure normalization.
They lie in \(H^\infty\).  The resulting sequence is the initial trace of
the actual differentiated history.
\end{proposition}

\begin{proof}
Begin with \(V_0(0)=u_0\).  Since \(H^\infty\) is closed under finite
products and spatial differentiation, and the Riesz transforms act
continuously on every \(H^m\), equation
\eqref{eq:pressure-recurrence} determines \(Q_0(0)\in H^\infty\).
Equation \eqref{eq:velocity-recurrence} then determines
\(V_1(0)\in H^\infty\).  If the rows through height \(n\) have been
determined in \(H^\infty\), the finite sum in
\eqref{eq:pressure-recurrence} determines \(Q_n(0)\), and
\eqref{eq:velocity-recurrence} determines \(V_{n+1}(0)\).
Induction gives the complete sequence.  On every strict classical
subinterval, direct differentiation of \eqref{eq:ns} obeys the same
recurrence with the same base value, so the recursively generated sequence
is the trace of the actual jet.
\end{proof}

\begin{remark}[Temporal response and material response]
The symbols \(V_1=u_t\) and \(V_2=u_{tt}\) denote time derivatives at fixed
spatial position.  If \(D_t:=\partial_t+u\cdot\nabla\), then the material
acceleration is \(D_tu\), not \(V_1\); its next material response is
\(D_t^2u\), not \(V_2\).  The proof uses the fixed-position jet because it
is the jet generated by repeated differentiation of the evolution equation.
\end{remark}

\section{Finite compatible rectangles and the complete intersection}

Write \(\Istar=(0,T_*)\).  For \(N\in\Nzero\), \(M\ge1\), and
\(0<T<T_*\), define
\begin{equation}
\begin{aligned}
\frakR_{N,M}(u;T)
:={}&\sum_{n=0}^{N}
\norm{\partial_t^n u}_{L^2(0,T;H^{M+1})}^2 \\
&+\sum_{n=0}^{N}
\norm{\partial_t^{n+1}u}_{L^2(0,T;H^{M-1})}^2.
\end{aligned}
\label{eq:rectangle-norm}
\end{equation}
To type the records precisely, let \(\mathscr V_{N,M}\) be the space of
tuples \(\mathbf v=(v_0,\ldots,v_{N+1})\) of divergence-free distributions
for which
\begin{equation}
\norm{\mathbf v}_{\mathscr V_{N,M}}^2
:=\sum_{n=0}^{N}\norm{v_n}_{L^2(\Istar;H^{M+1})}^2
+\sum_{n=0}^{N}\norm{v_{n+1}}_{L^2(\Istar;H^{M-1})}^2
<\infty.
\label{eq:velocity-record-space}
\end{equation}
Let
\[
\mathscr P_{N,M}
:=\prod_{n=0}^{N}L^1(\Istar;H^{M+1}(\R^3))
\]
with its finite product topology.  Because \(M+1>3/2\), the products below
belong to the stated pressure factors.  Define \(\mathscr X_{N,M}\) as the
subset of \(\mathscr V_{N,M}\times\mathscr P_{N,M}\), with the relative
product topology, whose records
\[
(\mathbf v,\mathbf q)
=\bigl((v_n)_{0\le n\le N+1},(q_n)_{0\le n\le N}\bigr)
\]
satisfy, for \(0\le n\le N\),
\begin{align}
\partial_t v_n&=v_{n+1},
\label{eq:record-time-link}\\
q_n&=R_iR_j\sum_{a=0}^{n}\binom{n}{a}v_{a,i}v_{n-a,j},
\label{eq:record-pressure-link}\\
v_{n+1}
&=\nu\Delta v_n
-\sum_{a=0}^{n}\binom{n}{a}(v_a\cdot\nabla)v_{n-a}
-\nabla q_n
\label{eq:record-equation-link}
\end{align}
in \(\mathcal D'(\Istar\times\R^3)\).  Thus
\[
\cR_{N,M}[u_0]
:=\bigl((V_n)_{0\le n\le N+1},(Q_n)_{0\le n\le N}\bigr)
\in\mathscr X_{N,M}
\]
denotes the datum-specific finite restriction of the actual jet, not the
ambient record space or another construction of the solution.

\Needspace{14\baselineskip}
\begin{theorem}[Datum-generated whole-terminal compatible intersection]
\label{thm:intersection}
For every \(N\in\Nzero\) and \(M\ge1\), the datum \(u_0\), the parameter
\(\nu\), and
their one maximal pressure-complete history generate a rectangle
\(\cR_{N,M}[u_0]\) satisfying
\begin{equation}
\boxed{
\sup_{0<T<T_*}\frakR_{N,M}(u;T)
\le C_{N,M}(u_0,\nu,T_*)<\infty.}
\label{eq:whole-terminal}
\end{equation}
The constant may depend on the finite coordinate \((N,M)\); no constant
uniform in all derivative orders is asserted.  If \(N'\le N\) and
\(1\le M'\le M\), the restriction maps satisfy
\begin{equation}
\rho_{N',M'}^{N,M}\!\left(\cR_{N,M}[u_0]\right)
=\cR_{N',M'}[u_0].
\label{eq:compatibility}
\end{equation}
Every shared coordinate is the corresponding derivative of the same
velocity--pressure history.
\end{theorem}

The supremum in \eqref{eq:whole-terminal} is the decisive quantifier.  A
bound that changes as \(T\uparrow T_*\) would not pay the endpoint passage.
The theorem supplies one datum-dependent bound through the entire alleged
terminal horizon for each requested finite coordinate.

\paragraph{Passage to the whole terminal interval.}
Every summand in \eqref{eq:rectangle-norm} is nonnegative and increases with
\(T\).  Monotone convergence and \eqref{eq:whole-terminal} give, for
\(0\le n\le N\),
\begin{equation}
\partial_t^n u\in L^2(\Istar;H^{M+1}),
\qquad
\partial_t^{n+1}u\in L^2(\Istar;H^{M-1}).
\label{eq:whole-terminal-rows}
\end{equation}
The symbols in \eqref{eq:whole-terminal-rows} still denote the
distributional derivatives of the original history.  No terminal
representative has yet been added.

If \(N'\le N\) and \(1\le M'\le M\), define
\[
\rho_{N',M'}^{N,M}:
\mathscr X_{N,M}\longrightarrow
\mathscr X_{N',M'}
\]
by discarding temporal rows above \(N'+1\) and pressure rows above \(N'\), applying
the continuous embeddings
\[
H^{M+1}\hookrightarrow H^{M'+1},
\qquad
H^{M-1}\hookrightarrow H^{M'-1},
\]
to the retained coordinates.  Equation \eqref{eq:record-pressure-link}
keeps their pressure normalization fixed.
These maps satisfy
\begin{equation}
\rho_{N,M}^{N,M}=\operatorname{id},
\qquad
\rho_{N'',M''}^{N',M'}\circ\rho_{N',M'}^{N,M}
=\rho_{N'',M''}^{N,M}
\label{eq:bonding-composition}
\end{equation}
whenever \((N'',M'')\le(N',M')\le(N,M)\).
Equation \eqref{eq:compatibility} says that the actual datum-generated
records form a compatible family for this directed system.

\begin{proposition}[Projective assembly]
\label{prop:projective-assembly}
Define
\begin{equation}
\cI[u_0]
:=\bigl(\cR_{N,M}[u_0]\bigr)_{N\in\Nzero,\,M\ge1}.
\label{eq:intersection}
\end{equation}
Then \(\cI[u_0]\) is the datum-specific element of the inverse-limit space
\begin{equation}
\mathscr X_\infty
:=\varprojlim_{N,M}\mathscr X_{N,M}
=
\left\{
(\mathbf R_{N,M})_{N,M}:
\begin{aligned}
&\mathbf R_{N,M}\in\mathscr X_{N,M},\\[-0.2em]
&\rho_{N',M'}^{N,M}\mathbf R_{N,M}
=\mathbf R_{N',M'}
\quad\text{whenever }(N',M')\le(N,M)
\end{aligned}
\right\}.
\label{eq:inverse-limit-definition}
\end{equation}
Its velocity component satisfies
\begin{equation}
\boxed{
u\in\bigcap_{r,m\in\Nzero}H^r(\Istar;H^m(\R^3)).}
\label{eq:all-orders-intersection}
\end{equation}
In particular, its zeroth temporal projection satisfies
\begin{equation}
u\in\bigcap_{m\in\Nzero}L^2(\Istar;H^m(\R^3)).
\label{eq:zero-row}
\end{equation}
\end{proposition}

\begin{proof}
The identity \eqref{eq:compatibility}, together with the composition law
\eqref{eq:bonding-composition}, places the family
\eqref{eq:intersection} in the inverse limit
\eqref{eq:inverse-limit-definition}.  Every shared coordinate is already
the same distributional derivative of \((u,p)\), so there are no
representatives to choose or reconcile.

Fix \(r,m\in\Nzero\).  Choose \(N\ge r\) and \(M\ge m+1\).  For
\(0\le n\le r\), the first membership in
\eqref{eq:whole-terminal-rows} gives
\[
\partial_t^n u\in L^2(\Istar;H^{M+1})
\hookrightarrow L^2(\Istar;H^m).
\]
These are the distributional time derivatives of the same \(u\), which is
exactly the definition of \(u\in H^r(\Istar;H^m)\).  Since \(r\) and \(m\)
were arbitrary, \eqref{eq:all-orders-intersection} follows.  Setting
\(r=0\) gives \eqref{eq:zero-row}.
\end{proof}

The order matters.  The complete compatible intersection
\eqref{eq:all-orders-intersection} is now present.  The critical and
classical estimates that follow are consequences of it, not inputs used to
construct it.

\section{The intersection becomes an estimate generator}

\begin{lemma}[Hilbert-valued time trace]
\label{lem:hilbert-time-trace}
Let \(X\) be a Hilbert space and \(0<T<\infty\).  If
\(f\in H^1(0,T;X)\), then \(f\) has a unique representative in
\(C([0,T];X)\), and
\[
\norm{f(t)-f(s)}_X
\le |t-s|^{1/2}\norm{f_t}_{L^2(0,T;X)}
\qquad(0\le s\le t\le T).
\]
\end{lemma}

\begin{proof}
The Bochner fundamental theorem of calculus gives
\(f(t)-f(s)=\int_s^t f_\tau(\tau)\dd\tau\) for the absolutely continuous
representative.  Cauchy--Schwarz gives the displayed estimate and hence
strong continuity.  Two continuous representatives equal almost everywhere
only if they agree everywhere.  This is the standard one-dimensional trace
result; see \cite{LionsMagenes1972}.
\end{proof}

\begin{proposition}[Universal finite-order projection]
\label{prop:universal-projection}
For every \(r,m\in\Nzero\),
\[
\partial_t^r u\in C([0,T_*];H^m).
\]
If \(k\in\Nzero\), \(2\le p\le\infty\), and
\(m>k+\frac32-\frac3p\), then
\(\nabla^k\partial_t^r u\in L^q(0,T_*;L^p)\) for every
\(1\le q\le\infty\), with the bounds
\eqref{eq:master-estimate} and \eqref{eq:master-estimate-infinity} below.
\end{proposition}

\begin{proof}
The intersection changes the estimate problem.  Fix any finite temporal
order \(r\) and spatial order \(m\).  From
\eqref{eq:all-orders-intersection},
\[
u\in H^{r+1}(0,T_*;H^m),
\qquad
\partial_t^r u\in H^1(0,T_*;H^m).
\]
The one-dimensional Hilbert-valued Sobolev trace theorem gives
\[
\partial_t^r u\in C([0,T_*];H^m).
\]
Hence
\begin{equation}
\boxed{
M_{r,m}:=
\sup_{0\le t\le T_*}\norm{\partial_t^r u(t)}_{H^m}<\infty}
\qquad (r,m\in\Nzero).
\label{eq:master-trace}
\end{equation}
Every finite temporal derivative is uniformly bounded in every fixed finite
spatial Sobolev norm through \(T_*\).

Now fix \(k\in\Nzero\), \(2\le p\le\infty\), and choose
\[
m>k+\frac32-\frac3p.
\]
Spatial Sobolev embedding gives, for every \(t\in[0,T_*]\),
\[
\norm{\nabla^k\partial_t^r u(t)}_{L^p}
\le C_{k,m,p}\norm{\partial_t^r u(t)}_{H^m}.
\]
Combining this estimate with \eqref{eq:master-trace} yields, for
\(1\le q<\infty\),
\begin{equation}
\boxed{
\norm{\nabla^k\partial_t^r u}_{L^q(0,T_*;L^p)}
\le C_{k,m,p}\,T_*^{1/q}M_{r,m}<\infty,}
\label{eq:master-estimate}
\end{equation}
while
\begin{equation}
\norm{\nabla^k\partial_t^r u}_{L^\infty(0,T_*;L^p)}
\le C_{k,m,p}M_{r,m}<\infty.
\label{eq:master-estimate-infinity}
\end{equation}
\end{proof}

Equations \eqref{eq:master-estimate}--\eqref{eq:master-estimate-infinity}
are the projection rule.  A fixed
finite-order estimate names \((r,k,p,q)\); one selects a sufficiently high
finite Sobolev coordinate \(m\); the complete intersection supplies that
coordinate.  The next section carries this rule through the principal
classical and scale-sensitive readouts one by one.

\section{The adjacent critical block and its classical projections}

\begin{lemma}[Adjacent Sobolev-scale trace]
\label{lem:adjacent-trace}
Let \(s\in\R\).  If
\[
v\in L^2(\Istar;H^{s+1}),
\qquad
v_t\in L^2(\Istar;H^{s-1}),
\]
then \(v\) has a unique representative in \(C([0,T_*];H^s)\).  For
\(0\le a\le b\le T_*\),
\begin{equation}
\left|
\norm{v(b)}_{H^s}^{2}-\norm{v(a)}_{H^s}^{2}
\right|
\le
2\norm{v_t}_{L^2(a,b;H^{s-1})}
\norm{v}_{L^2(a,b;H^{s+1})}.
\label{eq:adjacent-trace-bound}
\end{equation}
\end{lemma}

\begin{proof}
Let \(A=(1-\Delta)^{1/2}\).  Time mollification and spatial Fourier
truncation first give, for smooth approximants,
\[
\norm{v(b)}_{H^s}^{2}-\norm{v(a)}_{H^s}^{2}
=2\operatorname{Re}\int_a^b
\pair{A^{s-1}v_t(t)}{A^{s+1}v(t)}_{L^2}\dd t.
\]
Cauchy--Schwarz gives \eqref{eq:adjacent-trace-bound}.  The assumed
\(L^2\) memberships let the approximants converge in the right-hand
pairing.  The limit has weak \(H^s\) continuity and a continuous \(H^s\)
norm; together these give strong \(H^s\) continuity.  Uniqueness follows
as in Lemma~\ref{lem:hilbert-time-trace}.  This is the adjacent-scale form
of the Lions--Magenes trace argument \cite{LionsMagenes1972}.
\end{proof}

\subsection{The critical adjacent block}

The first critical readout needs only one finite rectangle, but that
rectangle is selected after the full intersection has been assembled.
The coordinate \(\cR_{0,2}[u_0]\) gives
\[
u\in L^2(\Istar;H^3),
\qquad
u_t\in L^2(\Istar;H^1).
\]
In particular,
\begin{equation}
u\in L^2(\Istar;H^{5/2}),
\qquad
u_t\in L^2(\Istar;H^{1/2}).
\label{eq:critical-adjacent-pair}
\end{equation}
Lemma~\ref{lem:adjacent-trace} with \(s=3/2\) gives
\begin{equation}
u\in C([0,T_*];H^{3/2}).
\label{eq:critical-hthreehalf-trace}
\end{equation}

Let \(\Lambda=(-\Delta)^{1/2}\), and for \(s\ge0\) set
\[
E_s(t):=\frac12\norm{\Lambda^s u(t)}_2^2,
\qquad
H(t):=E_{1/2}(t).
\]
The pair \eqref{eq:critical-adjacent-pair}, justified first by time
mollification and frequency truncation, gives
\begin{equation}
\frac{\dd}{\dd t}E_{3/2}(t)
=\operatorname{Re}
\pair{\Lambda^{1/2}u_t(t)}{\Lambda^{5/2}u(t)}_{L^2}
\in L^1(\Istar).
\label{eq:ethreehalf-derivative}
\end{equation}
The zeroth row gives \(E_{3/2}\in L^1(\Istar)\).  Hence
\begin{equation}
E_{3/2}\in W^{1,1}(\Istar)\cap L^1(\Istar)
\subset L^\infty(\Istar).
\label{eq:ethreehalf-control}
\end{equation}
The stronger integer rows also give
\begin{equation}
H'(t)
=\operatorname{Re}
\pair{\Lambda^{1/2}u(t)}{\Lambda^{1/2}u_t(t)}_{L^2}
\in L^1(\Istar).
\label{eq:height-derivative}
\end{equation}

To retain the pressure response in the critical balance, define
\[
\mathcal P_{1/2}(t)
:=-\operatorname{Re}
\pair{\Lambda^{1/2}u(t)}
{\Lambda^{1/2}\bigl((u\cdot\nabla)u+\nabla p\bigr)(t)}_{L^2}.
\]
Pairing \eqref{eq:ns} with \(\Lambda u\) gives the exact identity
\begin{equation}
H'=-2\nu E_{3/2}+\mathcal P_{1/2}.
\label{eq:critical-balance}
\end{equation}
Equations \eqref{eq:ethreehalf-control},
\eqref{eq:height-derivative}, and \eqref{eq:critical-balance} yield
\begin{equation}
\boxed{
H\in W^{1,1}(\Istar),\qquad
E_{3/2}\in L^\infty(\Istar)\cap L^1(\Istar),\qquad
\mathcal P_{1/2}\in L^1(\Istar).}
\label{eq:critical-block}
\end{equation}
The critical height, the adjacent energy, and the pressure-complete
production current all reach the terminal face with exactly the time
regularity claimed in \eqref{eq:critical-block}.  No value at the endpoint
is assigned to the \(L^1\) current.

\subsection{Uniform Sobolev and higher-response control}

Taking \(r=0\) in \eqref{eq:master-trace} gives, for every finite \(m\),
\begin{equation}
\sup_{0\le t\le T_*}\norm{u(t)}_{H^m}<\infty.
\label{eq:uniform-sobolev}
\end{equation}
Any continuation interface based on the divergence of one fixed Sobolev norm
is excluded by \eqref{eq:uniform-sobolev}.  If \(m>5/2\), then
\(H^m(\R^3)\hookrightarrow W^{1,\infty}(\R^3)\), so
\[
\sup_{0\le t\le T_*}\norm{\nabla u(t)}_{L^\infty}<\infty.
\]
Choosing a higher finite row gives, for each fixed \(k\),
\[
\sup_{0\le t\le T_*}\norm{\nabla^k u(t)}_{L^\infty}<\infty.
\]
The temporal index in \eqref{eq:master-trace} is also arbitrary.  Thus, for
every fixed \(r,k\),
\begin{equation}
\sup_{0\le t\le T_*}
\norm{\nabla^k\partial_t^r u(t)}_{L^\infty}<\infty.
\label{eq:all-fixed-responses}
\end{equation}
The velocity and its entire finite temporal-response tower remain finite in
every requested finite spatial norm.

\subsection{Serrin estimates}

One classical continuation family asks for
\[
u\in L^q(0,T_*;L^p),
\qquad
\frac2q+\frac3p\le1,
\qquad p>3.
\]
This is the Prodi--Serrin scale \cite{Prodi1959,Serrin1962,Serrin1963}.
The intersection contains the concrete row
\[
u\in L^2(0,T_*;H^2(\R^3)).
\]
Since \(H^2(\R^3)\hookrightarrow L^\infty(\R^3)\),
\begin{equation}
\boxed{u\in L^2(0,T_*;L^\infty(\R^3)),}
\label{eq:serrin}
\end{equation}
and its indices satisfy
\[
\frac22+\frac3\infty=1.
\]
The single finite row \eqref{eq:serrin} already lies on the Serrin boundary.
More generally, \eqref{eq:master-estimate} supplies every fixed admissible
pair in the stated range by selecting a suitable spatial coordinate.

\subsection{The critical
\texorpdfstring{\(L_t^\infty L_x^3\)}{L-infinity in time, L3 in space}
endpoint}

The coordinate \((r,m)=(0,1)\) in \eqref{eq:master-trace} gives
\[
u\in C([0,T_*];H^1(\R^3)).
\]
The embedding \(H^1\hookrightarrow L^6\), together with interpolation
between \(L^2\) and \(L^6\), yields
\[
\norm{u(t)}_{L^3}
\le \norm{u(t)}_{L^2}^{1/2}\norm{u(t)}_{L^6}^{1/2}
\le C\norm{u(t)}_{H^1}.
\]
Consequently,
\begin{equation}
\boxed{u\in L^\infty(0,T_*;L^3(\R^3)).}
\label{eq:critical-l3}
\end{equation}
The critical endpoint that would have to fail at a finite singular time is a
finite coordinate projection of \(\cI[u_0]\).  Its role as an endpoint
regularity class is the theorem of Escauriaza, Seregin, and \v{S}ver\'ak
\cite{EscauriazaSereginSverak2003}; here the membership is obtained directly
from the compatible intersection.

\subsection{Vorticity, strain, and Lipschitz action}

Define
\[
\vort:=\nabla\times u,
\qquad
\strain:=\frac12\bigl(\nabla u+(\nabla u)^{\mathsf T}\bigr).
\]
The intersection contains
\[
u\in L^2(0,T_*;H^3).
\]
It follows that
\[
\vort,\strain,\nabla u
\in L^2(0,T_*;H^2)
\hookrightarrow L^2(0,T_*;L^\infty).
\]
Cauchy--Schwarz on the finite terminal interval gives
\begin{align*}
\int_0^{T_*}\norm{\vort(t)}_{L^\infty}\dd t
&\le T_*^{1/2}\norm{\vort}_{L_t^2L_x^\infty}<\infty,\\
\int_0^{T_*}\norm{\strain(t)}_{L^\infty}\dd t
&\le T_*^{1/2}\norm{\strain}_{L_t^2L_x^\infty}<\infty,\\
\int_0^{T_*}\norm{\nabla u(t)}_{L^\infty}\dd t
&\le T_*^{1/2}\norm{\nabla u}_{L_t^2L_x^\infty}<\infty.
\end{align*}
Thus the vorticity, strain, and Lipschitz actions all remain finite through
the alleged terminal time.  In the vorticity equation, the stretching term
is \((\vort\cdot\nabla)u\); its pointwise operator size is bounded by
\(\norm{\nabla u}_{L^\infty}|\vort|\).  The last integral above is thus the
accumulated Lipschitz action capable of amplifying vorticity, while the
strain integral isolates its symmetric stretching component.

\subsection{Critical height and the finite octave count}

Under the parabolic rescaling
\[
u^{(\lambda)}(x,t):=\lambda u(\lambda x,\lambda^2t),
\qquad
p^{(\lambda)}(x,t):=\lambda^2p(\lambda x,\lambda^2t),
\]
the homogeneous half-derivative norm is unchanged:
\[
\norm{\Lambda^{1/2}u^{(\lambda)}(t)}_2
=\norm{\Lambda^{1/2}u(\lambda^2t)}_2.
\]
This is the sense in which the height \(H=E_{1/2}\), defined before
\eqref{eq:critical-block}, is critical.  That block gives an absolutely continuous
representative on the closed terminal interval, and hence
\begin{equation}
H_{\max}:=\sup_{0\le t\le T_*}H(t)<\infty.
\label{eq:height-bound}
\end{equation}
Fix a positive reference height \(H_0\), and set \(H_k:=2^kH_0\).  The
\(k\)-th rung is reached when the same history attains the level \(H_k\).
If a doubling ladder reaches the levels
\[
H_0,\ 2H_0,\ 4H_0,\ldots,2^kH_0,
\]
then \(2^kH_0\le H_{\max}\).  Whenever \(H_0\le H_{\max}\),
\begin{equation}
k\le
\left\lfloor\log_2\frac{H_{\max}}{H_0}\right\rfloor.
\label{eq:octave-count}
\end{equation}
If \(H_0>H_{\max}\), no rung is reached.  An infinite terminal Zeno ladder
would be a sequence \(t_k\uparrow T_*\) with \(H(t_k)=H_k\) for every \(k\).
The finite count above excludes such a sequence: the intersection bounds
the critical height before the counting argument begins.

\subsection{Dyadic shell summations}

Let \((\Delta_j)_{j\ge-1}\) be an inhomogeneous Littlewood--Paley
decomposition as in \cite{BahouriCheminDanchin2011}.  For \(j\ge0\),
Bernstein's inequality gives
\[
\norm{\Delta_j\vort(t)}_{L^\infty}
\lesssim 2^{3j/2}\norm{\Delta_j\vort(t)}_{L^2}
=2^{-j/2}\,2^{2j}\norm{\Delta_j\vort(t)}_{L^2}.
\]
After handling the low block \(j=-1\) once, Cauchy--Schwarz in \(j\) yields
\begin{align*}
\sum_{j\ge-1}\norm{\Delta_j\vort(t)}_{L^\infty}
&\lesssim
\left(\sum_{j\ge0}2^{-j}\right)^{1/2}
\left(\sum_{j\ge0}2^{4j}
\norm{\Delta_j\vort(t)}_{L^2}^2\right)^{1/2}
+C\norm{\Delta_{-1}\vort(t)}_{L^2}\\
&\lesssim \norm{\vort(t)}_{H^2}
\lesssim \norm{u(t)}_{H^3}.
\end{align*}
Integrating in time and using the \(L_t^2H_x^3\) coordinate gives
\begin{equation}
\boxed{
\sum_{j\ge-1}\int_0^{T_*}
\norm{\Delta_j\vort(t)}_{L^\infty}\dd t
\lesssim T_*^{1/2}\norm{u}_{L_t^2H_x^3}<\infty.}
\label{eq:vorticity-shells}
\end{equation}
Tonelli's theorem permits the interchange of the nonnegative shell sum and
the time integral in \eqref{eq:vorticity-shells}.

The same selection mechanism pays any fixed finite shell weight.  Fix
\(\sigma\in\R\) and \(k\in\Nzero\), and choose an integer
\[
M>\sigma+k+\frac32.
\]
For \(j\ge0\),
\[
2^{\sigma j}\norm{\Delta_j\nabla^k u(t)}_{L^\infty}
\lesssim
2^{-(M-\sigma-k-3/2)j}
\,2^{Mj}\norm{\Delta_j u(t)}_{L^2}.
\]
The first factor is square summable in \(j\), while the second is controlled
by \(\norm{u(t)}_{H^M}\).  The low block is again finite.  Hence
\begin{equation}
\int_0^{T_*}\sum_{j\ge-1}2^{\sigma j}
\norm{\Delta_j\nabla^k u(t)}_{L^\infty}\dd t
\le C_{\sigma,k,M}T_*^{1/2}
\norm{u}_{L^2(0,T_*;H^M)}
<\infty.
\label{eq:weighted-shells}
\end{equation}
Each fixed finite weighted shell sum is controlled by one sufficiently high
finite Sobolev row.  A further sum over every Sobolev order would ask for
growth information across the tower, which is a different assertion.

\subsection{Pressure projections}

The pressure recurrence \eqref{eq:pressure-recurrence} can be written
directly as
\begin{equation}
\partial_t^r p
=R_iR_j\sum_{a=0}^{r}\binom{r}{a}
(\partial_t^a u_i)(\partial_t^{r-a}u_j).
\label{eq:pressure-row}
\end{equation}
Fix finite \(r,m\), and choose a finite spatial level high enough that all
factors on the right of \eqref{eq:pressure-row} lie continuously in a
Sobolev algebra.  Multiplication and the Riesz transforms are continuous at
that level.  More explicitly, the choice \(m+2\) gives
\begin{equation}
\norm{\partial_t^r p(t)}_{H^m}
\le C_{r,m}\sum_{a=0}^{r}
\norm{\partial_t^a u(t)}_{H^{m+2}}
\norm{\partial_t^{r-a}u(t)}_{H^{m+2}}.
\label{eq:pressure-product-bound}
\end{equation}
Each factor on the right is continuous in time by
\eqref{eq:master-trace}.  The master trace estimate then gives
\begin{equation}
\partial_t^r p\in C([0,T_*];H^m)
\qquad\text{for every fixed finite }r,m.
\label{eq:pressure-all-orders}
\end{equation}
Fix \(k\in\Nzero\), \(2\le p\le\infty\), and choose
\(m>k+\frac32-\frac3p\).  Applying the same Sobolev embedding used in
Proposition~\ref{prop:universal-projection}, followed by the finite time
embedding, gives
\begin{equation}
\nabla^k\partial_t^r p
\in L^q(0,T_*;L^p)
\qquad
\left(
r,k\in\Nzero,\ 2\le p\le\infty,\ 1\le q\le\infty
\right).
\label{eq:pressure-mixed-norms}
\end{equation}
For example, choosing \(m>5/2\) after one finite derivative margin gives
\[
\nabla p\in L^\infty(0,T_*;L^\infty)
\subset L^1(0,T_*;L^\infty).
\]
Pressure-based finite-order regularity interfaces are thus projections of
the same intersection, and pressure remains tied to the same velocity field
at all finite orders.

\subsection{Type-I terminal products}

Under the rescaling used above, with the terminal time rescaled along with
the solution, the quantities
\[
(T_*-t)^{1/2}\norm{u(t)}_{L^\infty},
\qquad
(T_*-t)\norm{\nabla u(t)}_{L^\infty}
\]
are dimensionless.  These two quantities are what the phrase
\emph{Type-I terminal products} means in this subsection; the label invokes no additional
continuation theorem.

From \eqref{eq:all-fixed-responses},
\[
\sup_{0\le t\le T_*}\norm{u(t)}_{L^\infty}<\infty,
\qquad
\sup_{0\le t\le T_*}\norm{\nabla u(t)}_{L^\infty}<\infty.
\]
Multiplication by the vanishing terminal factors gives
\begin{equation}
(T_*-t)^{1/2}\norm{u(t)}_{L^\infty}\longrightarrow0,
\qquad
(T_*-t)\norm{\nabla u(t)}_{L^\infty}\longrightarrow0
\label{eq:type-one}
\end{equation}
as \(t\uparrow T_*\).  The Type-I terminal products defined above do not
approach a positive terminal obstruction; they vanish.

\subsection{The ultraviolet tail}

Use the unitary Fourier convention
\[
\widehat u(\xi)
=(2\pi)^{-3/2}\int_{\R^3}e^{-ix\cdot\xi}u(x)\dd x.
\]
For \(K\ge1\), define the critical high-frequency tail
\[
H_{>K}(t)
:=\frac12\int_{|\xi|>K}|\xi|\,|\widehat u(t,\xi)|^2\dd\xi.
\]
Fix any real \(m>1/2\), and choose an integer \(\ell\ge m\).  The uniform
\(H^\ell\) bound gives a uniform \(\dot H^m\) bound.  On \(|\xi|>K\),
\[
|\xi|\le K^{1-2m}|\xi|^{2m}.
\]
It follows that
\[
H_{>K}(t)
\le\frac12K^{1-2m}\norm{u(t)}_{\dot H^m}^2.
\]
The uniform Sobolev bound \eqref{eq:uniform-sobolev} now gives
\begin{equation}
\boxed{
\lim_{K\to\infty}\sup_{0\le t\le T_*}H_{>K}(t)=0.}
\label{eq:uv-tail}
\end{equation}
Critical mass cannot persist at unbounded frequency as the finite terminal
time is approached.

\subsection{Exact scope of the projection principle}

The preceding arguments cover each requested fixed finite temporal order,
spatial order, shell weight, and mixed norm with \(2\le p\le\infty\).  They
do not assert one constant uniform over all orders.  The dyadic statements
are finite-horizon statements; extending their time integrals to an infinite
horizon would require added decay.  Small-data thresholds, symmetry classes,
and geometric vorticity-direction hypotheses retain their independent
assumptions.  None of those qualifications weakens the claim proved here:
every displayed fixed finite-order continuation estimate is obtained from a
finite coordinate of the already completed intersection.

\section{The intersection supplies one common endpoint}

The catalogue shows how widely the intersection projects.  Closure needs a
single object: the endpoint of the original velocity field.  The zeroth
temporal projection \eqref{eq:zero-row} gives, for every finite \(m\),
\[
u\in L^2(0,T_*;H^{m+2}).
\]
Let \(\mathbb P\) denote the Leray projection.  Sobolev multiplication gives
\begin{equation}
\norm{\Pdiv(u\otimes u)}_{H^m}
\le C_m\norm{u}_{H^{m+2}}^2.
\label{eq:nonlinear-bound}
\end{equation}
Consequently,
\[
\norm{\Pdiv(u\otimes u)}_{L^1(0,T_*;H^m)}
\le C_m\norm{u}_{L^2(0,T_*;H^{m+2})}^2<\infty.
\]
The linear term satisfies
\[
\norm{\Delta u}_{L^1(0,T_*;H^m)}
\le T_*^{1/2}\norm{u}_{L^2(0,T_*;H^{m+2})}<\infty.
\]
Applying \(\mathbb P\) to \eqref{eq:ns} gives
\begin{equation}
u_t=\nu\Delta u-\Pdiv(u\otimes u).
\label{eq:projected-equation}
\end{equation}
Equations \eqref{eq:nonlinear-bound}--\eqref{eq:projected-equation} imply
\[
u_t\in L^1(0,T_*;H^m).
\]
Since \(T_*<\infty\) and \(u\in L^2(0,T_*;H^m)\), Cauchy--Schwarz also
gives \(u\in L^1(0,T_*;H^m)\).  Hence
\begin{equation}
u\in W^{1,1}(0,T_*;H^m)
\hookrightarrow C([0,T_*];H^m)
\qquad(m\in\Nzero).
\label{eq:w11-endpoint}
\end{equation}
The representative in \eqref{eq:w11-endpoint} agrees almost everywhere
with the \(H^1_tH^m_x\) representative from
\eqref{eq:master-trace}; continuity makes them identical everywhere.
If \(k>m\), the \(H^k\)-continuous representative, viewed through the
embedding \(H^k\hookrightarrow H^m\), agrees in the same way with the
\(H^m\)-representative.  Thus the endpoint values are not selected
independently.  They are compatible traces of the same velocity field and
define one common endpoint
\begin{equation}
\boxed{
u_*:=u(T_*)\in\bigcap_{m\in\Nzero}H^m(\R^3)=H^\infty(\R^3).}
\label{eq:common-endpoint}
\end{equation}
Because divergence is continuous from \(H^1\) to \(L^2\), the identities
\(\nabla\cdot u(t)=0\) for \(t<T_*\) pass to
\[
\nabla\cdot u_*=0.
\]
The \(H^0=L^2\) trace also gives finite kinetic energy.  Letting
\(t\uparrow T_*\) in \eqref{eq:energy-before-endpoint} yields the energy
identity at the terminal face.

\section{The equation generates the actual endpoint jet}

The endpoint is not only a family of spatial norms.  It carries the full
incoming pressure-compatible jet.  Begin with
\[
V_0^*:=u_*.
\]
For \(n\in\Nzero\), set
\begin{align}
Q_n^*
&:=R_iR_j\sum_{a=0}^{n}\binom{n}{a}V_{a,i}^*V_{n-a,j}^*,
\label{eq:endpoint-pressure}\\
V_{n+1}^*
&:=\nu\Delta V_n^*
-\sum_{a=0}^{n}\binom{n}{a}(V_a^*\cdot\nabla)V_{n-a}^*
-\nabla Q_n^*.
\label{eq:endpoint-velocity}
\end{align}
Because \(V_0^*\in H^\infty\), the first pressure row belongs to every finite
Sobolev space after choosing a finite product margin, and
\(V_1^*\in H^\infty\).  Induction through
\eqref{eq:endpoint-pressure}--\eqref{eq:endpoint-velocity} gives
\begin{equation}
V_n^*,Q_n^*\in H^\infty
\qquad\text{for every }n\in\Nzero.
\label{eq:endpoint-jet}
\end{equation}

This recursion is not a free formal jet.  We identify it with the incoming
history row by row.  For every \(n,m\in\Nzero\), define
\[
\widehat V_n
:=\lim_{t\uparrow T_*}\partial_t^n u(t)
\quad\text{in }H^m.
\]
The limits are independent of \(m\) under Sobolev restriction, so
\(\widehat V_n\in H^\infty\).  Equation
\eqref{eq:pressure-all-orders} similarly defines
\[
\widehat Q_n
:=\lim_{t\uparrow T_*}\partial_t^n p(t)
\quad\text{in every }H^m.
\]

We prove \((\widehat V_n,\widehat Q_n)=(V_n^*,Q_n^*)\) by induction.
The base is \(\widehat V_0=u_*=V_0^*\).  Suppose
\(\widehat V_a=V_a^*\) for \(0\le a\le n\).  Fix a target Sobolev order
\(m\) and take all incoming limits in a higher finite order at which
multiplication, differentiation, and the Riesz transforms in the \(n\)-th
row are continuous.  Passing \(t\uparrow T_*\) in
\eqref{eq:pressure-recurrence} gives
\[
\widehat Q_n
=R_iR_j\sum_{a=0}^{n}\binom{n}{a}
V_{a,i}^*V_{n-a,j}^*
=Q_n^*.
\]
Passing through \eqref{eq:velocity-recurrence} then gives
\[
\widehat V_{n+1}
=\nu\Delta V_n^*
-\sum_{a=0}^{n}\binom{n}{a}
(V_a^*\cdot\nabla)V_{n-a}^*
-\nabla Q_n^*
=V_{n+1}^*.
\]
The target order \(m\) was arbitrary, so the identities hold in
\(H^\infty\).  Induction identifies every velocity and pressure row.
The alleged terminal point carries the complete pressure-compatible jet of
the history that reached it.

\section{Same-history restart and global closure}

\begin{lemma}[Uniform local restart]
\label{lem:local-restart}
Fix \(s>5/2\), \(\nu>0\), and \(B<\infty\).  There is
\(\delta=\delta(\nu,s,B)>0\) such that, for every \(t_0\in\R\), every
divergence-free
\(a\in H^s(\R^3)\) with \(\norm{a}_{H^s}\le B\) generates a unique strong
pressure-normalized solution on \([t_0,t_0+\delta]\).  If
\(a\in H^\infty\), every higher Sobolev order persists on the same interval,
the solution is smooth, and its right temporal jet at \(t_0\) is generated
by \eqref{eq:endpoint-pressure}--\eqref{eq:endpoint-velocity}.
\end{lemma}

\begin{proof}
This is the standard local \(H^s\) theory for \(s>5/2\), with the existence
time bounded below in terms of the fixed parameter and the \(H^s\) norm;
see \cite{FujitaKato1964,Kato1984,Temam2001}.  The pressure is recovered by
\eqref{eq:pressure}, so uniqueness of the velocity also fixes the normalized
pressure.  Persistence follows by applying the differentiated energy
estimates on the lifespan of the same strong solution.  Repeated right
differentiation at the initial face gives the normalized recurrences.
\end{proof}

Choose \(s>5/2\) and set
\[
B:=1+\sup_{0\le t\le T_*}\norm{u(t)}_{H^s}<\infty.
\]
Let \(\delta=\delta(\nu,s,B)\) be given by
Lemma~\ref{lem:local-restart}.  Choose \(t_0<T_*\) with
\(T_*-t_0<\delta/2\), and launch the local solution \(w\) from
\(w(t_0)=u(t_0)\).  It exists past \(T_*\).  On
\([t_0,T_*)\), both \(w\) and \(u\) are strong solutions with the same datum
at \(t_0\); uniqueness gives
\[
w(t)=u(t)\qquad(t_0\le t<T_*).
\]
Continuity in \(H^s\) then gives \(w(T_*)=u_*\).

Apply the same lemma at \(T_*\) to \(u_*\), and call the resulting
pressure-normalized restart \((z,q)\).  The restriction of \(w\) to the
right of \(T_*\) and \(z\) have the
same initial datum, so uniqueness identifies them on their common interval.
The right jet of \((z,q)\) is generated by
\eqref{eq:endpoint-pressure}--\eqref{eq:endpoint-velocity}; hence it equals
\((V_n^*,Q_n^*)_{n\ge0}\).  Section~7 identified this sequence with the
incoming jet.  The piecewise pair
\[
(\widetilde u,\widetilde p)(t)
:=
\begin{cases}
(u,p)(t),&0\le t<T_*,\\
(z,q)(t),&T_*\le t<T_*+\delta,
\end{cases}
\]
has matching derivatives of every order at \(T_*\), solves the same
normalized system there, and is smooth across the join.  By the overlap
argument it is the same history generated by \(u_0\), now continued beyond
\(T_*\).

This contradicts the maximality of a finite \(T_*\).  Hence
\begin{equation}
\boxed{T_*=\infty.}
\label{eq:global-time}
\end{equation}
The energy identity on the left and right pieces agrees at the common
\(L^2\) trace \(u_*\), so \eqref{eq:global-energy} holds throughout the
continued history.

The complete argument now closes along one datum-preserving chain:
\begin{equation}
\boxed{\begin{gathered}
u_0\longrightarrow(u,p)
\longrightarrow\{J_{n,\alpha},\Pi_{n,\alpha}\}
\longrightarrow\{\cR_{N,M}[u_0]\}_{N,M}
\longrightarrow\cI[u_0]\\
\longrightarrow u(T_*)\in H^\infty
\longrightarrow\text{actual pressure-complete endpoint jet}
\longrightarrow\text{same-history restart}
\longrightarrow T_*=\infty.
\end{gathered}}
\label{eq:closing-chain}
\end{equation}
Every consequence estimate used between the intersection and the endpoint
was a finite projection of \(\cI[u_0]\), not another history and not a
producer for the intersection.  This proves Theorem~\ref{thm:global}.

\hfill\(\square\)

\appendix

\section{Finite-order analytic tools used in the proof}

This appendix records the standard facts used above and the finite margins
at which they are applied.  No estimate in the paper takes a sum or a
supremum over all Sobolev orders.

\subsection{Sobolev products, embeddings, and multipliers}

In three dimensions,
\begin{equation}
H^s(\R^3)\hookrightarrow W^{k,p}(\R^3)
\qquad\text{if}\qquad
s>k+\frac32-\frac3p,
\quad 2\le p\le\infty.
\label{eq:appendix-sobolev-embedding}
\end{equation}
If \(s>3/2\), then \(H^s(\R^3)\) is an algebra, and the usual tame product
estimate gives
\[
\norm{fg}_{H^m}
\le C_{m,s}\bigl(
\norm{f}_{H^m}\norm{g}_{H^s}
+\norm{f}_{H^s}\norm{g}_{H^m}
\bigr)
\]
whenever the displayed norms are finite.  In particular, for each integer
\(m\ge0\),
\begin{equation}
\norm{\mathbb P\nabla\cdot(v\otimes w)}_{H^m}
\le C_m\norm{v}_{H^{m+2}}\norm{w}_{H^{m+2}}.
\label{eq:appendix-product}
\end{equation}

The Riesz transforms \(R_j=\partial_j(-\Delta)^{-1/2}\) and the Leray
projection \(\mathbb P\) are order-zero Fourier multipliers.  They act
continuously on \(H^s(\R^3)\) for every \(s\in\R\).  These facts justify
the pressure estimates \eqref{eq:pressure-product-bound}, the projected
equation \eqref{eq:projected-equation}, and the passage to every finite
endpoint-jet row.

\subsection{Time traces and compatible representatives}

For a Hilbert space \(X\), the Bochner space \(H^1(0,T;X)\) embeds into
\(C([0,T];X)\).  Lemma~\ref{lem:hilbert-time-trace} gave the direct
estimate used here.  The adjacent-scale version,
\[
L^2(0,T;H^{s+1})\cap
W^{1,2}(0,T;H^{s-1})
\hookrightarrow C([0,T];H^s),
\]
was proved in Lemma~\ref{lem:adjacent-trace}.  Both statements select unique
continuous representatives.  This uniqueness is what makes the terminal
values obtained at different Sobolev orders one compatible endpoint rather
than a collection of unrelated limits.

\subsection{Inhomogeneous dyadic convention}

Choose smooth radial cutoffs \(\chi,\varphi\) with
\[
1=\chi(\xi)+\sum_{j\ge0}\varphi(2^{-j}\xi),
\]
and define
\[
\Delta_{-1}:=\chi(D),
\qquad
\Delta_j:=\varphi(2^{-j}D)\quad(j\ge0).
\]
For \(j\ge0\), Bernstein's estimate reads
\[
\norm{\Delta_j f}_{L^\infty}
\le C2^{3j/2}\norm{\Delta_j f}_{L^2}.
\]
The inhomogeneous square-function equivalence
\[
\norm{f}_{H^M}^2
\simeq
\norm{\Delta_{-1}f}_2^2
+\sum_{j\ge0}2^{2Mj}\norm{\Delta_jf}_2^2
\]
is the input used in \eqref{eq:vorticity-shells} and
\eqref{eq:weighted-shells}; see
\cite{BahouriCheminDanchin2011}.  The low block is retained once, and all
high-block weights are summed only after a fixed finite \(M\) has been
chosen.

\subsection{Constant dependence}

The constants \(C_{N,M}(u_0,\nu,T_*)\) belong to the settled
whole-terminal rectangle theorem and may vary with the selected finite
rectangle.  Constants in the projection estimates may also depend on their
displayed finite indices, on \(T_*\), and on the selected rectangle
constant.  The proof uses no analytic radius and no uniform growth law
across derivative order.

\begingroup
\small
\begin{thebibliography}{99}

\bibitem{Leray1934}
J.~Leray,
\newblock Sur le mouvement d'un liquide visqueux emplissant l'espace,
\newblock \emph{Acta Mathematica} \textbf{63} (1934), 193--248.

\bibitem{Prodi1959}
G.~Prodi,
\newblock Un teorema di unicit\`a per le equazioni di
{Navier{--}\hspace{0pt}Stokes},
\newblock \emph{Annali di Matematica Pura ed Applicata} \textbf{48}
(1959), 173--182.

\bibitem{Serrin1962}
J.~Serrin,
\newblock On the interior regularity of weak solutions of the
{Navier{--}\hspace{0pt}Stokes} equations,
\newblock \emph{Archive for Rational Mechanics and Analysis} \textbf{9}
(1962), 187--195.

\bibitem{Serrin1963}
J.~Serrin,
\newblock The initial value problem for the {Navier{--}\hspace{0pt}Stokes} equations,
\newblock in \emph{Nonlinear Problems}, University of Wisconsin Press,
Madison, 1963, 69--98.

\bibitem{FujitaKato1964}
H.~Fujita and T.~Kato,
\newblock On the {Navier{--}\hspace{0pt}Stokes} initial value problem. I,
\newblock \emph{Archive for Rational Mechanics and Analysis} \textbf{16}
(1964), 269--315.

\bibitem{Kato1984}
T.~Kato,
\newblock Strong \(L^p\)-solutions of the {Navier{--}\hspace{0pt}Stokes} equation in
\(\mathbb R^m\), with applications to weak solutions,
\newblock \emph{Mathematische Zeitschrift} \textbf{187} (1984), 471--480.

\bibitem{EscauriazaSereginSverak2003}
L.~Escauriaza, G.~A. Seregin, and V.~\v{S}ver\'ak,
\newblock \(L_{3,\infty}\)-solutions of the {Navier{--}\hspace{0pt}Stokes} equations
and backward uniqueness,
\newblock \emph{Russian Mathematical Surveys} \textbf{58} (2003),
211--250.

\bibitem{LionsMagenes1972}
J.-L.~Lions and E.~Magenes,
\newblock \emph{Non-Homogeneous Boundary Value Problems and Applications,
Vol. I},
\newblock Springer, 1972.

\bibitem{Temam2001}
R.~Temam,
\newblock \emph{{Navier{--}\hspace{0pt}Stokes} Equations: Theory and Numerical
Analysis},
\newblock AMS Chelsea, 2001.

\bibitem{BahouriCheminDanchin2011}
H.~Bahouri, J.-Y.~Chemin, and R.~Danchin,
\newblock \emph{Fourier Analysis and Nonlinear Partial Differential
Equations},
\newblock Springer, 2011.

\end{thebibliography}
\endgroup

\end{document}
